\section{\ENG{System Context} και Θέση του \ENG{SCORM Engine} στο \ENG{LMS} Οικοσύστημα}
\label{sec:system_context_scorm}

Η υλοποίηση που εξετάζεται δεν συγκροτεί ένα ενιαίο \ENG{LMS} με εσωτερικό \ENG{SCORM} υποσύστημα, αλλά ένα μικρό σύνολο συνεργαζόμενων υπηρεσιών με διακριτά όρια ευθύνης. Στο επίπεδο της αλληλεπίδρασης του χρήστη, το σημείο εισόδου παραμένει το \ENG{Laravel LMS} του αποθετηρίου \texttt{\ENG{example-lms-client}}. Εκεί υλοποιούνται οι διαχειριστικές ενέργειες εισαγωγής μαθημάτων, συγχρονισμού χρηστών και εγγραφών, καθώς και οι σελίδες του εκπαιδευόμενου για εκκίνηση μαθήματος και προβολή προόδου. Το \ENG{LMS}, ωστόσο, δεν εκτελεί απευθείας \ENG{SCORM} περιεχόμενο και δεν ενσωματώνει δικό του \ENG{run-time API adapter}. Η ευθύνη αυτή μεταφέρεται σε εξωτερικά υποσυστήματα.

Ο \ENG{SCORM Engine} του αποθετηρίου \texttt{\ENG{scorm-engine}} αποτελεί τον κεντρικό \ENG{backend} πυρήνα του οικοσυστήματος ο οποίος εκθέτει το βασικό \ENG{REST API} του συστήματος. Ο πυρήνας αυτός αναλαμβάνει την εισαγωγή και καταχώριση των πακέτων μαθημάτων, την ανάλυση του \texttt{\ENG{imsmanifest.xml}} στο απολύτως αναγκαίο επίπεδο για την εκτέλεση, τη διαχείριση των \ENG{launches}, την παραλαβή των \ENG{runtime commits} και τη μεταφορά της κανονικοποιημένης προόδου σε μόνιμες εγγραφές προσπαθειών.

Ο \ENG{player} του αποθετηρίου \texttt{\ENG{player}} τοποθετείται ανάμεσα στο περιεχόμενο και στον \ENG{browser}, ως ξεχωριστό \ENG{Node/TypeScript service}. Ο \ENG{player} αποδίδει τη σελίδα εκκίνησης, ζητά από το \ENG{engine} το \ENG{launch context}, κατεβάζει το συμπιεσμένο πακέτο από το \ENG{object storage}, το αποσυμπιέζει στον τοπικό του \ENG{cache} και εκθέτει τα αρχεία του μαθήματος μέσω του \ENG{route} \texttt{\ENG{/content/:launchId/*}}. Παράλληλα εισάγει στον \ENG{browser} ελαφρές \ENG{SCORM API} ως \texttt{\ENG{window.API}} ή \texttt{\ENG{window.API\_1484\_11}}, ώστε το \ENG{SCO} να αλληλεπιδρά με τη συνεδρία χωρίς άμεση πρόσβαση στα εσωτερικά \ENG{endpoints} του \ENG{engine}.

Ανάμεσα στο \ENG{LMS} και στον \ENG{engine} παρεμβάλλεται το \ENG{PHP SDK} του αποθετηρίου \texttt{\ENG{scorm-engine-php-sdk}}. Η ύπαρξη τυποποιημένων \texttt{\ENG{Api Service}} υποστηρίζουν το \ENG{integration layer}. Αυτό είναι κρίσιμο και σε αρχιτεκτονικό επίπεδο, διότι διαχωρίζει το \ENG{application-specific LMS logic} από το \ENG{transport contract} του \ENG{SCORM} υποσυστήματος.

Η υποδομή που υποστηρίζει τα παραπάνω περιγράφεται από το \texttt{\ENG{central-docker-infrastructure}}. Η \ENG{Postgres} χρησιμοποιείται ως κύρια σχεσιακή βάση για \ENG{courses}, \ENG{users}, \ENG{enrollments}, \ENG{attempts}, \ENG{launches} και στιγμιότυπα κατάστασης εκτέλεσης. Η \ENG{Redis} λειτουργεί ως προσωρινός χώρος αποθηκεύσης για τις ενεργές καταστάσεις εκτέλεσης των \ENG{launches}. Το \ENG{MinIO} εξυπηρετεί την αποθήκευση των \ENG{ZIP} πακέτων μαθημάτων, ενώ το \ENG{ELK stack} υποστηρίζει τη συγκέντρωση και επιθεώρηση \ENG{logs}.

\begin{figure}[H]
  \centering
\begin{chapterfigurebox}
  \begin{tikzpicture}[
      font=\footnotesize,
      node distance=1.9cm and 2.4cm,
      box/.style={draw, rounded corners, fill=gray!5, align=center, text width=3.2cm, minimum height=1.1cm},
      emph/.style={draw, rounded corners, fill=gray!12, align=center, text width=3.5cm, minimum height=1.1cm, font=\footnotesize\bfseries},
      store/.style={draw, cylinder, shape border rotate=90, aspect=0.28, minimum height=1.2cm, minimum width=3.2cm, align=center, fill=gray!5},
      arrow/.style={-Latex, thick},
      boundary/.style={draw, dashed, rounded corners, inner sep=11pt}
    ]
    \node[box] (browser) {\ENG{Browser}\\\ENG{Learner / Admin}};
    \node[box, below=of browser] (lms) {\ENG{Laravel LMS}};
    \node[box, right=of lms] (sdk) {\ENG{PHP SDK}};
    \node[emph, right=of sdk] (engine) {\ENG{SCORM Engine}};
    \node[emph, above=of engine] (player) {\ENG{SCORM Player}};

    \node[store, below left=2.0cm and 0.5cm of engine] (postgres) {\ENG{Postgres}\\\ENG{courses / attempts}};
    \node[store, below=2.0cm of engine] (redis) {\ENG{Redis}\\ενεργή κατάσταση εκτέλεσης};
    \node[store, below right=2.0cm and 0.5cm of engine] (minio) {\ENG{MinIO}\\\ENG{course ZIP objects}};
    \node[box, right=2.4cm of player] (elk) {\ENG{ELK Stack}\\\ENG{logs and inspection}};

    \draw[arrow] (browser) -- (lms);
    \draw[arrow] (lms) -- (sdk);
    \draw[arrow] (sdk) -- (engine);
    \draw[arrow] (browser.east) to[bend left=18] (player.west);
    \draw[arrow] (player) -- (engine);
    \draw[arrow] (engine) -- (postgres);
    \draw[arrow] (engine) -- (redis);
    \draw[arrow] (engine) -- (minio);
    \draw[arrow] (player) -- (minio);
    \draw[arrow] (engine) -- (elk);
    \draw[arrow] (player) -- (elk);

    \node[boundary, fit=(lms)(sdk)] (lmsboundary) {};
    \node[font=\footnotesize, above=0.12cm of lmsboundary] {\ENG{Business / integration layer}};

    \node[boundary, fit=(engine)(player)(postgres)(redis)(minio)(elk)] (runtimeboundary) {};
    \node[font=\footnotesize, above=0.12cm of runtimeboundary] {\ENG{SCORM platform services}};
  \end{tikzpicture}
\end{chapterfigurebox}
  \caption{Πραγματικό \ENG{system context} της πλατφόρμας, με διακριτούς ρόλους για \ENG{LMS}, \ENG{SDK}, \ENG{engine}, \ENG{player} και υποστηρικτικά αποθετήρια.}
  \label{fig:platform_context}
\end{figure}

